\documentclass[12pt,a4paper]{article}

%COMPILING: IMPORTANT INFORMATION
% You'll need to generate a .bbl file from
% the .bib file first, otherwise the publications
% will not show up in the cv.

% in TeXstudio, this means
% 1. F1 or F6 for usual compiling (generate .aux)
% 1. F8 for bibliography compiling (generate .bbl)
% 2. F1 or F6 for usual compiling once more(generate .pdf)

% use your favourite citation style here
\usepackage{apacite}

\usepackage[left=2cm,right=2cm,top=2cm,bottom=2cm]{geometry}

\newcommand{\headingfont}{\Large\color{OliveGreen}}

\usepackage[dvipsnames]{xcolor}



\newenvironment{list1}{
  \begin{list}{\ding{113}}{%
      \setlength{\itemsep}{0in}
      \setlength{\parsep}{0in} \setlength{\parskip}{0in}
      \setlength{\topsep}{0in} \setlength{\partopsep}{0in}
      \setlength{\leftmargin}{0.17in}}}{\end{list}}
\newenvironment{list2}{
  \begin{list}{$\bullet$}{%
      \setlength{\itemsep}{0in}
      \setlength{\parsep}{0in} \setlength{\parskip}{0in}
      \setlength{\topsep}{0in} \setlength{\partopsep}{0in}
      \setlength{\leftmargin}{0.2in}}}{\end{list}}



% as suggested by the following two links:
%https://tex.stackexchange.com/questions/227933/using-bibentry-to-cite-in-text-reference-with-apacite
% and
%https://tex.stackexchange.com/questions/65348/clash-between-bibentry-and-hyperref-with-bibstyle-elsart-harv/65401#65401
\usepackage{bibentry}
\makeatletter\let\saved@bibitem\@bibitem\makeatother
\usepackage[colorlinks=true]{hyperref}
\makeatletter\let\@bibitem\saved@bibitem\makeatother

\interlinepenalty=10000 

% hanging publications: 1st line starts at beginning of line, 
% further lines are placed a bit to the right
\usepackage{hanging}
\newcommand\publication[1]{%
    \smallskip\par\hangpara{1.5em}{1}\bibentry{#1}\smallskip
}

\usepackage{array,longtable,amsmath}

\renewcommand{\baselinestretch}{1.2}


\begin{document}


\begin{center}
{\Huge Haim Bar, Ph.D.}\\
Associate Professor\\Department of Statistics\\
University of Connecticut\\\vspace{5mm}
haim.bar@uconn.edu   $\;\boldsymbol{\cdot}\;$ https://haim-bar.uconn.edu $\;\boldsymbol{\cdot}\;$  https://github.com/haimbar 
\end{center}

\section*{\sc{\headingfont Research Interests}}
Statistical models; shrinkage estimation; Bayesian high-dimensional data; high throughput applications in
biology; % (e.g., genomics, brain imaging); applications in econometrics; 
Bayesian statistics; machine learning.


\section*{\sc{\headingfont  Education}}
{\bf Cornell University}, Ithaca, New York USA.  Statistics: Ph.D. 2012, M.Sc. 2010.\\
{\bf Yale University}, New Haven, Connecticut USA. M.Sc. in Computer Science.\\
{\bf The Hebrew University}, Jerusalem, Israel. B.Sc. in Mathematics, Cum Laude.

\section*{\sc{\headingfont  Academic Experience}}

\noindent{\bf University of Connecticut}, Storrs, Connecticut, USA.
\begin{list1}
\item[]  2019 -- Present: Associate Professor.
\item[] 2013 -- 2019: Assistant Professor. 
\begin{list2}
\item[] I am an associate professor at the University of Connecticut, in the Department of Statistics. I am affiliated with
the Institute for Systems Genomics (ISG),  the Institute for Collaboration on Health, Intervention, and Policy (InCHIP),
and the CT Institute for the Brain and Cognitive Sciences (IBACS).
\item[] Teaching: STAT 5005: Introduction to Applied Statistics;  
STAT 5505: Applied Statistics I;
STAT 5605: Applied Statistics II;
STAT 3515Q/5515: Design of Experiments; % (\textit{Excellence in teaching recognition - Spring 2014, Spring 2017});
STAT 3675Q: Statistical Computing; % (\textit{Excellence in teaching recognition - Spring 2016});
STAT/BIST 5225: Data Management in SAS and R; % (\textit{Excellence in teaching recognition - Fall 2017});
STAT/BIST 5215: Statistical Consulting; % (\textit{Excellence in teaching recognition - Spring 2018, Fall 2018});
STAT 5615: Categorical Data Analysis.
\end{list2}
\end{list1}

\noindent{\bf The Hebrew University}, Jerusalem, Israel.
2022 (May-August). Visiting faculty.


\noindent{\bf Cornell University}, Ithaca, New York, USA.
\begin{list1}
\item[]  2011 -- 2013:  Statistical Consultant.
\begin{list2}
\item[] As a member of the Cornell Statistical Consulting Unit (CSCU) I provided
statistical consulting to Cornell faculty and staff.  I handled a diverse range
of applications, including economics, business, genomics, nutrition, sociology,
etc. I presented a wide range of topics in workshops, and as a guest lecturer in
courses at Cornell, including on Bayesian statistics, missing data, multiple
testing, etc.
\item[] Teaching: BTRY 6020: Statistical Methods II.
\end{list2}
\item[] 2012 -- 2013: PostDoc, Department of Statistical Science
\begin{list2}
\item[] I collaborated with the Melnick Lab at the Weill Cornell Medical College.
I developed new methods and applied existing
techniques to analyze very large data sets involving up to millions of
features.  Primarily, we focused on methylation data for epigenetic
analysis, and RNA sequencing for transcription regulation analysis.
\end{list2}
\end{list1}


%\vspace{-.3cm}
%{\em Research Assistant} \hfill {\bf July 2011 - December, 2011}
%\begin{list1}
%\item[] I worked with Professor Francesca Molinari on the theory and application
%of sharp identification regions.  Implemented a fast and accurate estimation
%method based on Support Vector Machines.   We applied the method to linear
%prediction, and game theory problems.
%\end{list1}
%
%\vspace{-.3cm}
%{\em Graduate Student} \hfill {\bf August 2005 - December 2011}
%\begin{list1}
%\item[] M.Sc. and Ph.D.~ programs.
%\end{list1}
%
%\vspace{-.3cm}
%{\em Head Teaching Assistant} \hfill {\bf August 2009 - December 2010}
%\begin{list1}
%\item[] BTRY 6010, Statistical Methods I (\textit{Outstanding TA
%Award} 2010-2011).
%%Responsibilities included preparing lab and homework material; administrative
%%management; and supervising four teaching assistants.
%\end{list1}
%
%
%\vspace{-.3cm}
%{\em Teaching Assistant} \hfill {\bf August, 2005 - May 2010}
%\begin{list1}
%\item[] BTRY 6020, Statistical Methods II (graduate level class).
%\item[] BTRY 6010, Statistical Methods I (graduate level class).
%\item[] ORIE 474, Statistical Data Mining.
%\item[] CS 101M, Introduction to Computer Programming
%(\textit{Outstanding TA Award}, 2006).
%\item[] MATH 335/COMS 480, An Introduction to Cryptology.
%\end{list1}



    % bibliography
    \nobibliography{hbrefs.bib}
    \bibliographystyle{apacite} 

    \section*{\sc{\headingfont Publications}}

    % optional: split into different subsections for
    % peer-reviewed articles, conference abstracts, ...

    % example with links
%        (\href{https://arxiv.org/abs/1706.06969}{link}, \href{https://arxiv.org/pdf/1706.06969.pdf}{pdf}, \href{https://github.com/rgeirhos/object-recognition}{data and materials})

    % example without links
\publication{MCAP26}
\publication{NAR2025}
\publication{Bar15052025}
\publication{Frontiers2025}
\publication{kane_2025}
\publication{bar_ting_2025}
\publication{HMM2024} % add volume, issue, number
\publication{BMCBioinformatics2024}
\publication{axioms13090641}
\publication{ConvGeom2023}
\publication{bar_yan_2023}
\publication{Witherow2023}
\publication{Fumagalli2023}
\publication{BARRETT2022}
\publication{KATNENI2022}
\publication{nutrients2021}
\publication{QREM}
\publication{GenMed2021}
\publication{AJHG2021}
\publication{JDS2021}
\publication{PlosOne2021}
\publication{Environmetrics2021}
\publication{SEMMS}
\publication{SLAS2020}
\publication{THROMRES2020}
\publication{JMB2020}
\publication{F1000Research}
\publication{ijms2019}
\publication{SciRep2019}
\publication{STAT2019}
\publication{AEM2019}
\publication{MGGM2019}
\publication{JMB2019}
\publication{CHM2019}
\publication{WICS2019}
\publication{AJPGP2019}
\publication{ASMBI2018}
\publication{CDN2018}
\publication{JASA2017}
\publication{CE2017}
\publication{EPP2017}
\publication{ArchToxicol2017}
\publication{Haemophilia2017}
\publication{Nutrients2016}
\publication{MHR2016}
\publication{ICSA2004}
\publication{JASA2014}
\publication{RepSci2014}
\publication{PlosOne2014}
\publication{PLANTCELL2014}
\publication{JFE2014}
\publication{BLOOD2013}
\publication{SIM2012}
\publication{SAGMB2012}
\publication{PlosOne2012b}
\publication{FASEB2012}
\publication{PlosOne2012}
\publication{StatMod2011}
\publication{FASEB2011}
\publication{JEB2011}
\publication{STS2010}
\publication{BioresTech2011}
\publication{IEEE2003}

    % the cumbersome way
%   \hangpara{1.5em}{1}Knuth, D. E. (1998). 
%   \textit{The art of computer programming: sorting and searching} (Vol. 3). Pearson Education.


\section*{\sc{\headingfont Under review or in preparation}}

\begin{list1}
%\item[-] Bar, H. \textit{Interprocess Communication – How to Talk to Your Programs}. \textbf{Revise and resubmit}.
%\item[-] Bar, H.; Booth, J.; Wells, M.T. \textit{Automated Model Building and Goodness-of-fit via Quantile Regression}.
%\item[] Liu, K.; Bar, H.  \textit{Large-P Variable Selection in Two-Stage Models}.
\item[-] Nguyen, H.,  Bar, H.,  Chi, Z,  Pozdnyakov, V. \textit{A Time-Varying Branching Process Approach to Model Self-Renewing Cells}. 
\item[-] Bar, H., Nguyen, H., Conover, J. \textit{A State Model for the Analysis of Stem Cell
Proliferation and Differentiation}
%\item[-] Jamie
%\item[-] Assaf
\end{list1}
  

\section*{\sc{\headingfont Grants and Contracts}}
\begin{list1}
\item[-] Research Foundation of The City University of New York, (May 1, 2024 - April 30, 2027). \textit{Interaction of Choline and Fat in the Prenatal Programming of Nonalcoholic Steatohepatitis}, Funding amount: \$113,536 .
\item[-] CLAS Research in Academic Themes 2024 funding initiative (January 2024-June 2024). PI. \textit{The Blessing of Dimensionality - Theory of Convex Geometry and Concentration of Measure with Applications in Machine Learning, Finance, and Genomics}. Funding amount \$30,000.
\item[-] BlackRock Inc. additional funding for \textit{The Blessing of Dimensionality - Theory of Convex Geometry and Concentration of Measure with Applications in Machine Learning, Finance, and Genomics}. Funding amount \$15,000.

\item[-] CLAS Research in Academic Themes 2022 funding initiative (January 2023-June 2024). PI. \textit{A novel statistical approach to uncover the hidden language of long non-coding RNAs}. Funding amount \$49,725.
\item[-] IBACS Seed Grant. PI. \textit{Modeling and visualizing the formation of brain cavities covering from stem cells.} Funding period: 1/1/2022-12/31/2022. Funding amount \$18,516.
\item[-] NIH Award: \textit{Characterization of the ventricular-subventricular stem cell niche during human brain development}. 
PI: J. Conover. H. Bar (co-Investigator.) Funding period: April 2020-March 2024. %Funding amount \$2,010,500.
\item[-] NSF Award  1612625. PI. \textit{Variable Selection in the High Dimensional, Low Sample Size Setting - Beyond the Linear Regression and Normal Errors Model.}
Funding period: 8/15/16-7/31/19. Funding amount \$150,000.
\item[-] National Science Center - Poland.
\textit{Identification of transcriptomic markers of maize resistance to cereal aphids}.
PI Dr. Hubert Sytykieticz.
Funding period: 2/9/17-2/8/20.
\item[-] Travelers.
\textit{Modeling and analysis of large insurance claim and occurrence data: A partnership between UConn and Travelers}. 2018 -- PI: Dr. Dipak Dey, co-PIs: Haim Bar, Kun Chen, Elizabeth Schifano, Xiaojing Wang.
Funding period: 8/1/17-7/31/18. Funding amount \$137,881.
2019 -- PI: Dey, D. Co-Investigators: Haim Bar, Kun Chen,  Victor
Hugo Lachos Davilla. Funding period: 8/1/2018 - 7/31/2019. Funding amount \$146,451.
 \end{list1}


\section*{\sc{\headingfont Patents}}
\begin{list1}
\item[-] U.S. Patent 20230041627 A1 (Publication Date: February 9, 2023)\\
{\em Assay Compound Screening.}
\item[-] U.S. Patent 9384677 (granted July 5, 2016)\\
U.S. Patent 20150213730 A1 (Publication Date: July 30, 2015)\\
U.S. Patent 9076342 B2 (granted July 7, 2015)\\
U.S. Patent 20090208910 A1 (Publication Date: August 20, 2009)\\
{\em Automated execution and evaluation of network-based training exercises.}\\
Inventors: Brueckner, S. K.; Adelstein, F. N.; Bar, H.; Donovan, M.
\item[-] U.S. Patent 20150143355 (Publication Date: May 21, 2015)\\
{\em Service Oriented Architecture Version and Dependency Control.}\\
Inventors: Tingstrom, D. J.; Joyce, R. A.; Stillerman, M. A.; Brueckner, S. K.; Bar, H.
\item[-] U.S. Patent 8,984,396 (granted March 17, 2015)\\
{\em Identifying and representing changes between extensible markup language (XML) files using symbols with data element indication and direction indication.}\\
Inventors: Tingstrom, D. J.; Joyce, R. A.; Stillerman, M. A.; Brueckner, S. K.; Bar, H.
\item[-] U.S. Patent 8,898,285 (granted November 25, 2014)\\
{\em Service oriented architecture version and dependency control}\\
Inventors: Tingstrom, D. J.; Joyce, R. A.; Stillerman, M. A.; Brueckner, S. K.; Bar, H.
\item[-]  U.S. Patent 8,286,249 (granted October 9, 2012)\\
U.S. Patent 7,748,040 (granted June 29, 2010);\\
{\em Attack correlation using marked information.}\\
Inventors: Adelstein F. N., Bar H., Alla P. and Proskourine N.
\end{list1}

\section*{\sc{\headingfont Short Courses}}
\begin{list1}
\item[-] Data Science. University of Connecticut, Pre-College Summer, 2022-2024.
\item[-] Advanced statistical modeling ideas (session on missing data). The University of Barcelona, Spain. January, 2019.
\item[-] Advanced topics in R and Machine Learning. The Technion, Haifa, Israel. July, 2018.
\item[-] R/Bioconductor workshop - the Technion, Haifa. November, 2016 (co-taught with Dr. Martin Morgan).
\end{list1}


\section*{\sc{\headingfont Invited Talks}}
\begin{list1}
\item[]\textit{High Dimensional Space Oddity.} 
\begin{list2}
\item[-] The Joint Statistical Modeling, Nashville, TN (August, 2025).
\item[-] Statistics Department Colloquium, UConn (October, 2024).
\end{list2}
\item[]\textit{On graphical models and convex geometry.}
\begin{list2}
\item[-] Consortium for Data Analytics in Risk, Berkeley, CA. November, 2023.
\item[-] SIAM Conference on Financial Mathematics and Engineering. Philadelphia, PA. June, 2023.
\item[-] Mathematics and Statistics Department, Old Dominion University. March, 2023.
\item[-] Statistics Department, The Hebrew University, Jerusalem, Israel. June 2022.
\end{list2}
\item[] \textit{Large-P Variable Selection in Two-Stage Models.}
\begin{list2}
\item[-] CMStatistics, London, UK,  December 2020.
\item[-] University of Connecticut, March 2020.
\item[-] University of Haifa, Israel. May, 2020.
\end{list2}
\item[] \textit{Quantile Regression Modelling via Location and
Scale Mixtures of Normal Distributions.}
\begin{list2}
\item[-] Cornell University, October 2019.
\end{list2}
\item[] \textit{A Mixture Model to Detect Edges in Sparse Co-expression Graphs.}
\begin{list2}
\item[-] The 3rd Eastern Asia Meeting on Bayesian Statistics, Seoul, South Korea. July, 2018.
\item[-] Yonsei University, Seoul, South Korea. July, 2018.
\item[-] Korea University, Seoul, South Korea. July, 2018.
\item[-] Bayesian Modeling, Computation, and Applications, Storrs, CT, May 2018.
\end{list2}
\textit{A Scalable Empirical Bayes Approach to Variable
Selection in Generalized Linear Models.}
\begin{list2}
\item[-] The 34th Quality and Productivity Research Conference, Storrs, CT, June 2017.
\item[-] The 10th ICSA International Conference, Shanghai, China. December, 2016.
\item[-] Temple University, Research Colloquium, Department of Statistical Science, April 2017.
\item[-] Baruch College, CUNY, Colloquium, Department of Information Systems and Statistics Zicklin School of Business, November 2016.
\end{list2}
\item[]\textit{A Scalable Empirical Bayes Approach to Variable Selection.}
\begin{list2}
\item[-] 2016 Joint Statistical Meetings, Chicago, IL. July, 2016.
\item[-] Joint UCONN/UMASS Statistics Colloquium, University of Massachusetts Amherst, MA. October 7, 2015
\item[-] Statistics Seminar, Department of Mathematics - University of Maryland. October 1, 2015.
\item[-] Colloquium - Mathematical Biosciences Institute - The Ohio State University.
Sep. 21, 2015
\end{list2}
\item[]\textit{An Empirical Bayes Approach to Variable Selection and QTL Analysis.}
\begin{list2}
\item[-] Purdue University, Research Colloquium, Statistics Department, October 2014.
\item[-] Modern Modeling Methods ($M^3$) Conference, Storrs, CT, May 2014.
\end{list2}
\item[]\textit{Model-based approaches for big-data problems, with applications in genomics.}
\begin{list2}
\item[-] University of Connecticut, Institute for Systems Genomics Annual Networking Workshop. Storrs, CT. May 13, 2014.
\end{list2}
\textit{A Bivariate Model for Simultaneous Testing in Bioinformatics Data}
\begin{list2}
\item[-] 3rd International Conference and Exhibition on Biometrics \& Biostatistics, Baltimore, MD, October 2014.
\item[-] University of Iowa, February 2013
\item[-] NIH/NCI, March 2013
\item[-] University of Connecticut, February 2013
\item[-] University of Rochester, February 2013
\item[-] Cornell University, October 2012
\end{list2}
\begin{list2}
\item[-]\textit{Accounting for Heaping in Retrospectively Reported Event Data - A
Mixture Model Approach.}
ICSA Applied Statistics Symposium in New York City, NY, USA. June 2011.
\end{list2}
\textit{A Heap of Trouble? Accounting for Mismatch Bias in Retrospectively
Collected Data on Smoking.}
\begin{list2}
\item[-] 3rd Biennial Conference of the American Society of Health Economists, Ithaca, NY,
June 2010.
\end{list2}
\end{list1}

\section*{\sc{\headingfont Conference Presentations}}
\begin{list1}
\item[-] \textit{An Empirical Bayes Approach to Variable Selection and QTL Analysis}. Frontiers Of Hierarchical Modeling In Observational Studies, Complex Surveys And Big Data: A Conference Honoring Professor Malay Ghosh. College Park, MD. May 29-31, 2014.
\item[-] \textit{A Mixture-Model Approach for Testing for Unequal Variances in Microarray
Data.}
Conference of Applied Statistics Ireland, Galway, Ireland, 2011.
\item[-] \textit{An Empirical Bayes Approach to Variable Selection and QTL Analysis.}
In the Proceedings of the 25th International Workshop on Statistical Modelling,
Glasgow, Scotland, 2010.
\end{list1}

\section*{\sc{\headingfont Journal Referee}}
Annals of Applied Statistics;
Bayesian Analysis;
Bioinformatics;
Biometrics;
Behavior Research Methods;
BMC bioinformatics;
Conservation Biology;
Evaluation and Program Planning;
International Journal of Environmental Research and Public Health;
International Statistical Review;
Journal of Agricultural, Biological, and Environmental Statistics;
Journal of Computational and Graphical Statistics;
Journal of Statistical Modeling;
Journal of Statistical Planning and Inference;
Journal of Statistical Theory and Practice;
Journal of the American Statistical Association;
Methodology \& Computing in Applied Probability;
PLOS ONE;
Statistical Analysis and Data Mining;
Statistical Modelling;
Statistics and Its Inference;
Statistics in Medicine;
Stats.

\section*{\sc{\headingfont Referee - Other}}
\begin{list1}
\item[-] National Science Foundation (NSF), CISE.
\item[-] National Science Foundation (NSF), 2017.
\item[-] The Donaghue Medical Research Foundation:\\
 \hspace*{5mm}``Greater Value Portfolio'' - Review Board Member, 2015, 2016, 2017.\\
 \hspace*{5mm}``Another Look'' - Review Board Member, 2016.
\item[-] National Aeronautics and Space Administration - Review Board Member, 2016.
\item[-] International Chinese Statistical Association - Student Award Committee for the 2016 ICSA International Conference in Shanghai, China, Dec. 19-22, 2016.
\item[-] New England Statistics Symposium - IBM Student Award Committee, 2015.
\end{list1}





%
\section*{\sc{\headingfont Conference Activities}}
\begin{list1}
\item[-] Organizer: The Blessings of Dimensionality Workshop, University of Connecticut, Storrs, CT, 2024.
\item[-] Organizing committee: The Department of Statistics 60th anniversary, University of Connecticut, Storrs, CT, 2022.
\item[-] Organizing committee: The 33rd New England Statistics Symposium, Hartford, CT, 2019.
\item[-] Co-chair of the organizing committee: The 31st New England Statistics Symposium, Storrs, CT, 2017.
\item[-] Co-chair of the organizing committee: The 34th Quality and Productivity Research Conference, Storrs, CT, 2017.
\item[-] Session chair - Joint Statistical Meeting, 2016 (Chicago, IL).
Session title: R Tools for Statistical Computing
\item[-] Session chair - Joint Statistical Meeting, 2015 (Seattle, WA).
Session title: Methods in Machine Learning and Data Mining.
\item[-] Session chair - 3rd International Conference and Exhibition on Biometrics \&
 Biostatistics, 2014.
\end{list1}

\section*{\sc{\headingfont Department and University Committees}}
\begin{list1}
\item[] Distinguished Statistician Colloquium - the Pfizer Series (Chair).
\item[] Graduate program curriculum - applied statistics (Chair).
\item[] Computing (Chair).
\item[] New England Statistics Society.
\item[] New England Statistics Symposium, 2017.
\item[] Gratis faculty appointment (Chair).
\item[] Visiting assistant professor search committee.
\item[] Graduate students and distinguished alumni awards.
\item[] UConn's Q-Center.
\item[] Environmental health and safety (Chair).
\item[] Library / tech reports (Chair).
\end{list1}

\section*{\sc{\headingfont Other Services}}
\begin{list1}
\item[] College of Liberal Arts and Sciences. Research Advisory Committee AY22-23.
\item[] The New England Journal of Statistics in Data Science - Software section. Editor.
\item[] Data Science in Science. Associate Editor.
\item[] Sankhya Series A, the Indian Journal of Statistics. (2018-2021) Associate Editor.
\item[] Health Policy Statistics Section communication officer (July 2016-June 2019).
\item[] Health Policy Statistics Section - best student paper review board (for JSM 2017).
\item[] The 10th ICSA International Conference - Young Researcher Award review committee.
\item[] New England Statistics Society - co-chair of the committee to establish the society's journal.
\item[] Eastern North American Region (ENAR) 2019: Program committee member. (2018).
\end{list1}



\section*{\sc{\headingfont Ph.D. Advising}}
\begin{list1}
\item[] Kangyan Liu, 2019 (major advisor)
\item[] Paul McLaughlin, 2019 (major advisor)
\item[] Cheng Huang, 2022 (associate advisor)
\item[] Srawan Bishoni, 2022 (associate advisor)
%\item[] Katherine Zavez (associate advisor)
\item[] Ziqi Yang, 2021  (associate advisor)
\item[] Ellis Shaffer, 2020 (associate advisor)
\item[] Zhe Wang, 2020  (associate advisor)
\item[] Dakota Cintron, 2020  (associate advisor)
\item[] Chaoran Hu, 2020  (associate advisor)
\item[] Yulia Sidi, 2020  (associate advisor)
\item[] Jieying Jiao, 2020  (associate advisor)
\item[] Renjie Chen, 2019 (associate advisor)
\item[] Ruochen Zha, 2018 (associate advisor)
\item[] chongliang Luo, 2017 (associate advisor)
\item[] Aditya Mishra, 2017 (associate advisor)
\item[] Neha Mishra, 2017 (associate advisor)
\item[] Chantal Larose, 2015 (associate advisor)
\item[] Guang Ouyang, 2015 (associate advisor)
\end{list1}


\section*{\sc{\headingfont Professional Experience}}
 %{\bf Sunbird Software Consulting}, Ithaca, NY
%
% \vspace{-.3cm}
% {\em Consultant} \hfill {\bf December 2007 - April 2011}\\
% Provided software and statistical consulting to FreePatentsOnline.com and
%Appinions.com.
% \\ FreePatentsOnline.com provides on-line users with access to a large
%database of patents.
% My responsibilities included the integration of the German patent data, and
%the development of a patent-family database (using the DOCDB database from the
%European Patent Office).
% \\ Appinions.com uses state-of-the-art natural language processing tools to
%extract sentiments from documents.  I worked on the O2O project (Opinions to
%Outcomes), where I designed and developed software that performed statistical
%analysis to detect significant changes in sentiment toward publicly traded
%companies, as reflected in on-line publications, such as blogs and news feeds.

\noindent{\bf ATC-NY}, Ithaca, New York, USA.
\begin{list1}
\item[]  2003 -- 2010:  Principal Scientist.
\begin{list2}
\item[] Involved in several government-funded research projects, mostly related to
secure and survivable computer systems.  For example, in the SPRINT project
(``Secure Programming Using Artificial Intelligence Techniques'') I designed and
prototyped a tool to support secure programming with SPARK (SRI's Procedural
Agent Realization Kit).
 In ``Policy Projector'', a Phase I OSD-funded effort, I was involved in the
design of a tool that rapidly reveals potential problems with network policy
changes prior to their implementation.  In ``Software Pedigree Analysis''
(SPAN), a DARPA Phase I SBIR effort, I developed a prototype that maintains
software pedigree to support future forensic investigations; and in ``Active
Smart Targets for Effective Response'' (ASTER), a DARPA-funded Phase II SBIR, I
developed a prototype to identify hackers by feeding them traceable information
during exploratory probes.
\end{list2}
\end{list1}


  
\noindent{\bf MicroPatent LLC}, East Haven, Connecticut USA.
\begin{list1}
\item[]  2002 -- 2003:  Director of Software Development.
\begin{list2}
\item[] Responsible for all software development projects, leading a team of 15-20
programmers.
Designed and implemented an extensive patent-family database with over 50
million records, merging several different sources of data. Developed analytical
tools for advanced patent searchers.
\end{list2}
\item[] 2000 -- 2002: Special Projects Manager.
\begin{list2}
\item[] Led a team of five people to develop a large-scale system to deliver patent
databases in XML format.  The project included designing a database, writing a
program to convert the data to XML, writing programs to validate the XML, and
programs to automate and monitor the delivery process.
\end{list2}
\item[] 1999 -- 2000: Team Leader
\begin{list2}
\item[] Led a team that designed and implemented comprehensive patent and trademarks
databases and web interfaces.  In addition to developing web-based applications,
I was involved in other projects, including the design and implementation of a
Client/Server application to provide customers with direct access to the patent
database (using OpenText's ``Bibliographic Retrieval Services (BRS)/Search'');
design and implementation of an SQL-based database to save user queries, and a
corresponding web interface that allows users to construct complicated searches
by combining results from previous searches, and share searches with colleagues;
and an information-theoretic approach to improve retrieval of records.
\end{list2}
\item[] 1997 -- 1999: Unix Programmer
\begin{list2}
\item[] Involved in almost all the development projects in the company, including
programs to enable our web customers to download patent data in PDF format (the
first web site in the industry to do that), and in Lotus Notes format.
\end{list2}
\end{list1}



\noindent{\bf Motorola Israel - Wireless Access Systems Division}, Tel Aviv, Israel.
\begin{list1}
\item[]  1995 -- 1997:  Software Engineer.
\begin{list2}
\item[] As part of the Software Quality Assurance (SQA) team in the ``Wireless Local
Loop'' (WiLL) project, I was involved in a wide range of activities, including
extensive test and integration of wireless and land telephony systems, from the
operating system of the base stations to the end-user wireless equipment;
software engineering practices and procedures; writing a graphic user interface
for the WiLL system; writing a program to automate the analysis of massive log
files in order to detect anomalies;
and integration and testing of an encryption device in the fixed wireless
terminals.
\end{list2}
\end{list1}



\section*{\sc{\headingfont Other Skills}}
Statistical Packages:  R, BUGS, Matlab, JMP, SAS, Stata, SPSS.\\
Programming languages:  Pyhton, C, Perl, Java, JavaScript, Unix shell scripts, SQL, PHP,
HTML/XML, Tcl/Tk.

\vspace{2in}
Updated on \today


\end{document}
